\documentclass[dvipdfmx]{jreport}
\DeclareFontShape{JY1}{mc}{m}{it}{<->ssub*mc/m/n}{}
\DeclareFontShape{JY1}{mc}{m}{sl}{<->ssub*mc/m/n}{}
\DeclareFontShape{JY1}{mc}{m}{sc}{<->ssub*mc/m/n}{}
\DeclareFontShape{JY1}{gt}{m}{it}{<->ssub*gt/m/n}{}
\DeclareFontShape{JY1}{gt}{m}{sl}{<->ssub*gt/m/n}{}
\DeclareFontShape{JY1}{mc}{bx}{it}{<->ssub*gt/m/n}{}
\DeclareFontShape{JY1}{mc}{bx}{sl}{<->ssub*gt/m/n}{}
\DeclareFontShape{JT1}{mc}{m}{it}{<->ssub*mc/m/n}{}
\DeclareFontShape{JT1}{mc}{m}{sl}{<->ssub*mc/m/n}{}
\DeclareFontShape{JT1}{mc}{m}{sc}{<->ssub*mc/m/n}{}
\DeclareFontShape{JT1}{gt}{m}{it}{<->ssub*gt/m/n}{}
\DeclareFontShape{JT1}{gt}{m}{sl}{<->ssub*gt/m/n}{}
\DeclareFontShape{JT1}{mc}{bx}{it}{<->ssub*gt/m/n}{}
\DeclareFontShape{JT1}{mc}{bx}{sl}{<->ssub*gt/m/n}{}
\usepackage{xcolor}
\usepackage{framed}
\usepackage{pifont}
\usepackage[dvipdfmx,hidelinks]{hyperref}
\usepackage{pxjahyper}
\usepackage{bm}
\usepackage{amsfonts}
\usepackage{amsmath}
\usepackage{amssymb}
\usepackage{amsthm}
\usepackage[english]{babel}
\newtheorem{theorem}{Theorem}[section]
\newtheorem{corollary}{Corollary}[theorem]
\newtheorem{lemma}[theorem]{Lemma}
\newtheorem{definition}{Definition}[section]
\newtheorem{claim}{Claim}
\begin{document}
\title{比較可能定理}
\author{趙永明}
\date{2022-02-15}
\maketitle
\newpage
\tableofcontents
\addcontentsline{toc}{chapter}{contents}
\newpage
\chapter{比較可能定理}
\label{cha:比較可能定理}
\begin{theorem}
    任意の整列集合$(A,R)$に対して， $(A,R)$と順序同型となる順序数$\alpha$が一意的に存在する．
\end{theorem}
\begin{proof}
    一意性は順序数の間の同型写像は恒等写像に限ることから導かれる．\\
    以下存在性のみを示す:\\
    まず$a \in A$に対し，$W\langle a \rangle := \{ x \in A \mid xRa\} $ を$A$の $a$による切片と定める:\\
    続いて，集合 $G:=\{a \in A \mid \exists !\alpha \in ON; W \langle a \rangle \simeq \alpha \}$とする．\\
	$a \mapsto \alpha$により定まる写像 $f$とする．\\
以下$\exists \gamma \in ON; f(G) = \gamma$であることを示す:
	$\alpha \in f(G)$を任意に取る．この時， $\forall \beta  < \alpha ; \beta \in f(G)$を示せばいい．\\
	特に， $W \langle a \rangle \simeq \alpha$に対し，その同型写像を$h_a$と定める:\\
    この時，$\beta < \alpha$を任意に取る， $h_a$が特に全射であるため，$\exists c \in W\langle a\rangle>; h_a(c) = \beta$

	この時，$h_c = h_a\upharpoonright W \langle c \rangle $ と定めると，
$h_c : W\langle c \rangle \to h_a(W\langle c \rangle)$が同型写像になる．

この時, $c \in G \land \beta \in f\left( G \right) $ が示された. Lemma 4.4 より,$f\left( G \right)  \in ON$ が得られる.

$G \neq A$の時,  $A\setminus G$ の最小元$e$ が取れて, $G = W \langle e \rangle$となり, $f(G) \in ON$より,  $e \in G$が得られ，矛盾が生じる.

\end{proof}

\end{document}
